\documentclass{abstract}

% Bibliography
\usepackage[
    backend=biber,
    style=authoryear,
    doi=true,
    sorting=none,
    sortcites=true,
    maxbibnames=2,
    minbibnames=1,
    maxcitenames=2,
    mincitenames=1,
    citestyle=numeric-comp,
    giveninits=true,
    isbn=false,
    date=year
]{biblatex}
\addbibresource{abstract.bib}

% For email linking
\usepackage{hyperref}

\title{ENFP429 NIST Abstract}

\author{Jacob Myers}
%I'm not really sure on this part, figured I'd take a stab
\keywords{Cone Calorimeter, Cone Calorimetry, NIST, LaTeX}
\begin{document}
\maketitle
\section{Abstract}
% Big picture first: comment on how the data is currently available, then note the core elements of what we're trying to create [describe the key elements of the cone DB in your own words: structure, core elements/how to interact, how this is stored data/metadata]\\

The overarching goal of this project is to take data from cone calorimetry tests performed by NIST/the NBS over the past 40 years, currently stored in a physical form, and create a digital, searchable database. The intent is that if someone has a material that they wish to know the data for, they can easily find tests performed on that material in the past, along with the relevant test conditions and results.  


% Once the big picture is clear, then you can state the critical steps involved (for an abstract, big picture first; details of exactly what will come in the 2 pager summary). I encourage you to look at the lines your wrote below and see you you can organize them in a series of clear, succinct, critical/main idea steps. This should make sense to someone who has never heard of this project before.
The first step of this process is digitization. The physical data sheets are scanned, and then AI tools are used to take the data from the scanned image and create a digital data file. This file is then parsed into two smaller files. One is referred to as the data file, which contains the data resulting from the tests, like mass and energy emission, and the other is the metadata file, which contains the details of the test, like what energy input level was used, what the starting mass and time were, and the orientation of the material. The next step is to check these data and metadata files for mistakes. The process of checking the data and metadata files against the original scanned documents and editing them to remove transcription errors is termed “SmURFing.” An additional part of SmURFing is creating a Material ID, a name for the files that contains the material and primary testing conditions so that anyone can easily find the test that best matches what they are looking for.


% Now you can highlight what specifically you're doing (again, key concepts matter most for an abstract)
My primary duty in my role is performing this SmURFing. I access data and metadata files for a given test, and check it against the information from the original scan. I edit it as necessary, and add a Material ID. I then upload this changed version to the central NIST repository for it to be manually checked over by NIST employees. My particular data has been on composite materials. 
% You may wish to then add a conclude / add brief line or two highlighting the types of Data you'v SmURFed so far, "cables, composites, woods.. "


%% This kind of summary is fine for a student abstract, but wouldn't typically be included in a technical/research abstract. Since you have it, let's keep it but I'll encourage you to write more technically/clearly so an external person could understand in 30 s exacly what new skills/knowledge you have. XYZ tools/software + fire dynamics concepts (for this part, in your 2 page summary, you may wish to include a few lines of what the cone is/how it works/what's the controlling mechanisms here that we can study as part of your intro to the project // demosntration of why this [and the cone in general] is important)

We have additionally been educated on the function of a cone calorimeter, as well as the type of data it collects, and how this data is usually represented. We have had additional instruction on how to identify patterns of mass loss curves in cone calorimeter data. We have learned how to work with GitHub, in connecting to a central repository, pulling it down, and uploading to it. LaTeX has been used, and we have learned how to utilize LaTeX to create a document. 

 \cite{babrauskas1982development, brown1988cone,emil1989flammability}


\section{References}
%apologies, I cannot remember why we commented out the references section and if I was supposed to have references here.
%\printbibliography[heading=none]

\end{document}